\chapter{El verdadero poder de los humanos}\label{sec:el-verdadero-poder-de-los-humanos}

Se suponía que los humanos eran cobardes.

El registro de especies de la Federación Galáctica los tenía listados con un 2 sobre 16 en el índice de agresividad. Nuestras interacciones con la Unión Terrana hasta este punto respaldaban esas conclusiones. No habían librado guerras entre sí en siglos y habían formado un gobierno mundial unificado antes de lograr los viajes superlumínicos. Respondieron con entusiasmo en lugar de hostilidad al primer contacto, a diferencia de muchas especies.

La Tierra había resuelto todas sus disputas a través de la diplomacia y el compromiso desde que se convirtió en miembro oficial de la Federación. Por ejemplo, hace unos años, los expansionistas xanik reclamaron una colonia minera terrana como su territorio. La Federación se preparó para un conflicto menor, ya que esperaban que los humanos defendieran su puesto avanzado. Pero los humanos simplemente se encogieron de hombros y aceptaron ceder el planeta a cambio de un pequeño canon anual. En lugar de ir a la guerra, los terranos terminaron siendo importantes socios comerciales de los xanik.

También hubo un incidente en el que los paranoicos hoda'al arrestaron a embajadores terranos bajo cargos de espionaje. Encarcelar a diplomáticos sin evidencia alguna fue una clara provocación a la guerra, pero los humanos no hicieron nada. ¡Ni siquiera atacaron la instalación donde se encontraban sus representantes! Simplemente abrieron negociaciones por vías secundarias con los hoda'al y organizaron un intercambio de prisioneros, canjeando a algunos contrabandistas por su gente.

Las opiniones sobre los humanos variaban según a quién le preguntaras. Algunos en la Federación encontraban su pacifismo encomiable y apreciaban su diplomacia equilibrada. Otros pensaban que era la debilidad lo que los llevaba a evitar la guerra. Yo estaba en el último grupo; la única razón para no responder a insultos evidentes con agresión era que no tenían el ingenio ni la fuerza para hacerlo.

Cuando llegaron los devoradores, las tres especies más militaristas de la galaxia (según el índice de agresividad) se unieron para enfrentar su avance. No sabíamos mucho sobre ellos, pero los llamamos devoradores porque su única misión era succionar la energía de las estrellas. No puedo decirles por qué harían algo así. Sea cual sea su motivo, tomaban un sistema por la fuerza, lo agotaban y pasaban al siguiente.

Nuestra flota, la mejor que la Federación tenía para ofrecer, sufrió grandes pérdidas cuando chocamos con los destructores enemigos. Luchamos con todas nuestras fuerzas, pero no importó. Nuestras armas apenas parecían hacerles un rasguño. Fue una decisión difícil, pero ordené que lo que quedaba de la flota se retirara. Por mucho que necesitáramos detenerlos, perderíamos toda la armada si nos quedábamos más tiempo.

Envié una señal de socorro, transmitiendo nuestra sombría situación y suplicando refuerzos. Había otras especies con fuerzas militares menores, pero aun así poderosas, dentro de la Federación. Pero mi solicitud fue recibida con silencio. Ni uno solo de esos cobardes se ofreció como voluntario para ayudar. Al enterarse de nuestra derrota, supongo que decidieron huir y defenderse por su cuenta.

Pensé que estábamos solos, hasta que detectamos naves humanas saltando a nuestra posición. Qué irónico: los únicos que vinieron en nuestra ayuda fueron los pusilánimes de la galaxia. Según nuestros sensores, solo había cinco de ellas, lo cual no era ni de cerca suficiente para librar una batalla. Un espectáculo patético, pero era más de lo que habían enviado las otras potencias de la Federación.

—Señor, los terranos nos están contactando. ¿Qué creen que van a hacer, matar al enemigo a base de palabras? —bromeó el primer oficial Blez.

Escuché algunas risas de mi tripulación, pero los callé rápidamente. —Necesitamos toda la ayuda que podamos conseguir. En pantalla.

Un humano de cabello oscuro apareció en la pantalla. —Naves de la Federación, aquí el comandante Mikhail Rykov, de la Unión Terrana. Estamos aquí para ayudar en todo lo posible.

Incliné la cabeza con gratitud. —Gracias por venir, comandante Rykov. Soy el general Kilon. Por favor, únanse a nuestra formación y ayúdennos a cubrir nuestra retirada.

 —¿Retirada? —el comandante humano parpadeó un par de veces, luciendo confundido—. Nuestra intención es enfrentar y aniquilar al enemigo.

—Con solo cinco naves? Con todo respeto, los devoradores se cuentan por miles y aplastaron nuestra flota de igual magnitud. No esperaría que una especie pacífica como la suya comprendiera la guerra, pero les conviene seguir nuestro liderazgo —dije.

El comandante Rykov parecía aún más confundido. —¿Creen que los humanos son una especie pacífica? ¿Qué demonios? ¿Por qué pensarían eso?

—Bueno... nunca pelean con nadie. Resuelven todo con palabras. Ustedes son la especie con la calificación más baja en el índice de agresividad —respondí.

—Ya veo. La Federación nos ha juzgado mal en ese aspecto. ¿Sabe por qué evitamos la guerra, general?

 —¿Porque no creen que pueden ganar? ¿Miedo?

El humano rió de buena gana. —No, es porque sabemos qué somos. De lo que somos capaces. Y nadie se lo ha ganado todavía.

La idea de que los terranos hicieran amenazas ominosas habría sido una broma para mí hasta ahora, pero algo en el tono de Rykov me dijo que creía lo que decía con total convicción. Esto era un claro caso de delirio surgido de la falta de experiencia en la guerra interestelar. Los devoradores dejarían en ridículo a los terranos y los castigarían por su exceso de confianza. Sin embargo, si el comandante realmente quería enviar a sus hombres a una masacre, no lo detendría.

—Si insisten en luchar, ciertamente no me opondré. Pero sepan que están solos; nosotros nos largamos de aquí. ¿Cuál es su plan? —pregunté.

—Trajimos una nanobomba que desarrollamos. Nunca la hemos usado antes, ya que en aproximadamente el cinco por ciento de las simulaciones no se detiene con entidades localizadas y consume toda la materia del universo —el comandante Rykov dijo esto de manera demasiado casual para mi gusto—. Pero las programamos para autodestruirse después de unos segundos, lo cual probablemente funcionará. Alférez Carter, dispare contra el enemigo en cinco segundos.

Mis ojos se abrieron de par en par, alarmados. —Espere, espere, acaba de decir que podría destruirlo todo...

El buque insignia terrano disparó un misil antes de que pudiera pronunciar otra palabra para detenerlos. Al principio, pensé que habían fallado. El proyectil atravesó la flota devoradora sin impactar en ninguna nave. Luego, estalló en la retaguardia de la formación y se desató el infierno.

El espacio mismo pareció estremecerse cuando una explosión arrasó con todo a su alrededor. La fuerza fue tan poderosa que nuestros sensores solo pudieron proporcionar un mensaje de error como medición. Al menos un tercio de la flota devoradora se vaporizó al instante, cuando una cantidad improbable de energía y calor los convirtió en una sopa metálica. No había forma de que los ocupantes de esas naves sobrevivieran a eso.

Las naves enemigas más alejadas del epicentro sobrevivieron a la explosión inicial, aunque muchas de ellas sufrieron daños graves. Pero una fuerza invisible pareció diseccionarlas lentamente; solo pude mirar con incredulidad mientras los poderosos cruceros se desintegraban trozo a trozo. Supongo que la bomba había liberado una nube de nanitos que atacaron la estructura de las naves a nivel molecular.

Los devoradores apenas supieron qué los había golpeado. Para cuando pensaron en contraatacar, ya no quedaba nada con qué hacerlo. Su arsenal se evaporó en cuestión de segundos y, sin duda, su personal sufrió el mismo destino. Donde antes había un ejército imparable, ahora solo quedaba espacio vacío.

Los humanos habían desatado una ola de destrucción sin rival en toda mi carrera militar con un solo misil. El horror recorrió mis venas al pensar que algún día podrían dirigir sus monstruosas armas contra la Federación. No había forma de defenderse contra tales creaciones diabólicas.

El índice de agresividad necesitaba una actualización. El tipo de especie que inventaría armas como esas no era un 2. Mirando a mi tripulación, vi reacciones atónitas y horrorizadas que reflejaban las mías. Si alguna vez se volvían hostiles, los humanos representaban una amenaza de nivel máximo. Podrían, más que probablemente, arrasar toda la galaxia sin despeinarse.

—Ahora que eso está resuelto. ¡Deberían habernos invitado a la fiesta desde el principio! —sonrió el comandante Rykov—. Le diré algo, general: la próxima vez que nos veamos, nos debe una cerveza.

Fruncí el ceño. Los humanos podrían pedir mucho más que una bebida si quisieran. —Sí, creo que podemos hacer eso.

El comandante Rykov terminó la llamada y vi cómo las naves terranas saltaban de nuevo al hiperespacio. Todavía estaba tratando de asimilar todo lo sucedido, preguntándome cómo iba a ponerlo en palabras para el reporte de combate. La Federación no tenía idea de quiénes eran realmente los terranos, pero me aseguraría de que lo supieran.

Y mientras repasaba los acontecimientos del día en mi mente, comprendí. Finalmente entendí por qué una especie tan poderosa evitaba la guerra.

Los humanos evitan la guerra porque les sería demasiado fácil ganarla.